\subsection{问题一的模型建立与求解}
\label{sec:problem1}

[先总览本问要解决的目标、核心困难、主模型、输入输出、求解路线与验证方式，使后续四个小节形成完整闭环。]

\subsubsection{模型的建立}

[从问题特定关系出发定义变量、目标函数、约束、机理或统计结构，解释模型各部分的实际含义与依据。只有后文会用“由公式（3-1）”“根据式（3-1）”等表述明确引用的公式才使用 \texttt{CumcmReferencedEquation}，并设置 \texttt{\string\label}、在后文用 \texttt{\string\eqref} 引用；其余行内、行间公式一律不编号。]

\subsubsection{模型的求解}

[说明求解算法为何适用，写明输入输出、关键步骤、参数来源、停止条件、实现环境和复现条件。]

\subsubsection{模型的计算结果}

[集中呈现可追溯到代码输出或结果文件的真实数值、方案、分类或图表，写明单位、条件和有效数字。]

\subsubsection{问题一结果验证与解释}

[选择与本问风险相匹配的检验方法，解释结果的现实含义、可信度、失败情形与适用边界，不把相关或预测结果夸大为因果或真实判断。]

\FloatBarrier
